% Original applied explanation; see Chapter04/SOURCES.md.
\subsubsection{নোটিশের গুরুত্ব ও দেখতে মোটা লেখার পার্থক্য}
\begin{flushleft}\small
\begin{tabularx}{\linewidth}{|p{26mm}|X|X|}\hline
\textbf{বিবেচনা} & \textbf{\texttt{strong}} & \textbf{\texttt{b}}\\\hline
অর্থ & বিষয়বস্তুর বিশেষ গুরুত্ব গঠনে প্রকাশ করে & সাধারণভাবে দৃষ্টি আকর্ষণের অংশ চিহ্নিত করে\\\hline
দেখতে কেমন & মোটা দেখা যেতে পারে & প্রায় একইভাবে মোটা দেখা যেতে পারে; অর্থ এক নয়\\\hline
\end{tabularx}
\end{flushleft}
\textbf{উদাহরণ:} ভর্তি-নোটিশে \textbf{আবেদনের শেষ দিন} বিশেষ গুরুত্বপূর্ণ। \texttt{strong} দিয়ে সেই গুরুত্ব গঠনে প্রকাশ করা যায়।

নতুন বিষয় শুরু করতে শিরোনাম এবং বিস্তারিত লিখতে অনুচ্ছেদ ব্যবহার করবে। ফাঁকা জায়গা বানাতে বারবার \texttt{br} ব্যবহার করলে ছোট পর্দায় বিন্যাস অনিয়মিত হতে পারে। আলাদা অনুচ্ছেদে \texttt{p}, আর চিঠির ঠিকানার মতো একই অংশে প্রয়োজনীয় লাইনবিরতিতে \texttt{br} বেছে নেবে।
